\NativeTitle{13}{第二单元素养评价（B）}{（相交与平行）}{40}{6}{一 & 二 & 三 & 四 & 五 & 六}{ &  &  &  &  & }
\NativeSlot{276}{66}{%
\ExamGroup{一、填空。{\normalfont\kaishu\fontsize{10.8bp}{14bp}\selectfont （每空1分，共17分）}}

% p013_s01

\ExamQuestion{1.}{5}%
如下图，互相平行的是\mbox{（\hspace{1.8em}）}，互相垂直的是\mbox{（\hspace{1.8em}）}。\par
\BookFigure{figures/p013_s01-01.tex}
}
\NativeSlot{390}{36}{%
% p013_s02

\ExamQuestion{2.}{5}%
直线\(m\)与直线\(n\)相交成直角时，这两条直线的位置关系是\mbox{（\hspace{6em}）}，记作\mbox{（\hspace{1.8em}）}，读作\mbox{（\hspace{7em}）}；这两条直线的交点叫作\mbox{（\hspace{1.8em}）}。\par
}
\NativeSlot{441}{103}{%
% p013_s03

\ExamQuestion{3.}{5}%
钟表的圆形表盘上，除了3时整，\mbox{（\hspace{1.8em}）}时整的时针和分针所在直线也互相垂直。\par


\par\noindent%
\ExamTile{.74}{%
% p013_s04

\ExamQuestion{4.}{5}%
过春节贴春联是我国的传统习俗。如图，对联横批贴歪了，爷爷说：只要让横批的下边与门框的上边\mbox{（\hspace{1.8em}）}，横批就正了。\par
}\hfill%
\ExamTile{.235}{%
\BookFigure{figures/p013_s04-01.tex}
}\par
% p013_s05

\ExamQuestion{5.}{5}%
教室的推拉玻璃窗，推拉轨道的上下两条边互相\mbox{（\hspace{1.8em}）}，玻璃窗的竖直边框和轨道边的位置关系是\mbox{（\hspace{6em}）}。\par
}
\NativeSlot{616}{106}{%
\par\noindent%
\ExamTile{.67}{%
% p013_s06

\ExamQuestion{6.}{5}%
如图，直线\(EG\)平行于直线\(FH\)，如果线段\(EF\)的长度是5厘米，那么线段\(GH\)的长度是\mbox{（\hspace{1.8em}）}，线段\(AD\)的长度是\mbox{（\hspace{1.8em}）}，它们的长度\mbox{（\hspace{1.8em}）}，原因是\mbox{（\hspace{1.8em}）}，线段\(AB\)与线段\(AC\)的长度一定\mbox{（\hspace{1.8em}）}5厘米（填“大于”“小于”或“等于”）。\par
}\hfill%
\ExamTile{.305}{%
\BookFigure{figures/p013_s06-01.tex}
}\par

\par\noindent%
\ExamTile{.74}{%
% p013_s07

\ExamQuestion{7.}{5}%
在同一平面内，有4条直线，其中\(a\parallel b\)，\allowbreak{}\(b\perp c\)，\allowbreak{}\(a\)与\(c\)的位置关系是\mbox{（\hspace{6em}）}，\allowbreak{}\(b\perp d\)，\allowbreak{}\(b\perp c\)，那么\(c\)与\(d\)的位置关系\par
}\hfill%
\ExamTile{.235}{%
\BookFigure{figures/p013_s07-01.tex}
}\par
}
